\chapter{Dolbeault Cohomology}\label{dolbeault}
\chapterstatus{complete}

\section{Introduction}\label{dolbeault:section-introduction}

Dolbeault cohomology is the cohomology of the $\bar\partial$-operator on forms of type $(p,q)$. By the Dolbeault theorem
(Theorem~\ref{dolbeault:theorem-dolbeault}) it computes the sheaf cohomology of the holomorphic $p$-forms, and on a compact complex
manifold harmonic theory for $\bar\partial$ makes it finite-dimensional and gives Serre duality
(Theorem~\ref{dolbeault:theorem-serre-duality}). The numbers $h^{p,q} = \dim H^q(X, \Omega^p_X)$ are the \emph{Hodge numbers}; for
compact K\"ahler manifolds they are the dimensions of the pieces of the Hodge decomposition (Chapter~\ref{hodge-decomposition}).
References: \cite[Chapter 0]{griffiths-harris}, \cite[Chapter 4]{voisin-hodge-1}, \cite[Chapters 2--4]{huybrechts-complex},
\cite{wells-differential-analysis}.

\noindent\textbf{Prerequisites.} Chapter~\ref{holomorphic-bundles} (Holomorphic Vector Bundles) and Chapter~\ref{elliptic-operators}
(Elliptic Operators on Compact Manifolds).

$X$ is a complex manifold of dimension $n$ and $E$ a holomorphic vector bundle on it.

\section{The Dolbeault theorem}\label{dolbeault:section-dolbeault-theorem}

\begin{definition}\label{dolbeault:definition-dolbeault}
Let $\mathcal{A}^{p,q}(E)$ be the sheaf of smooth $E$-valued $(p,q)$-forms and $\mathcal{A}^{p,q}(X, E)$ its global sections. In a holomorphic frame of $E$ set $\bar\partial_E(\sum_j\alpha_j e_j) = \sum_j\bar\partial\alpha_j\,e_j$;
this is independent of the frame because the transition functions are holomorphic, and $\bar\partial_E^2 = 0$. The \emph{Dolbeault cohomology} is
$H^{p,q}_{\bar\partial}(X, E) = H^q(\mathcal{A}^{p,\bullet}(X, E), \bar\partial_E)$; for $E$ trivial we write $H^{p,q}_{\bar\partial}(X)$.
\end{definition}

\begin{lemma}\label{dolbeault:lemma-dbar-poincare-pq}
Let $\Delta' \subseteq \Delta$ be polydiscs with $\overline{\Delta'} \subseteq \Delta$ and $\alpha$ a smooth $(p,q)$-form on $\Delta$ with $q \geq 1$ and $\bar\partial\alpha = 0$. Then $\alpha = \bar\partial\beta$ on $\Delta'$ for some
smooth $(p, q-1)$-form $\beta$.
\end{lemma}

\begin{proof}
Writing $\alpha = \sum_I dz^I \wedge \alpha_I$ with $\alpha_I$ of type $(0,q)$, the condition $\bar\partial\alpha = 0$ holds componentwise, since $\bar\partial(dz^I) = 0$; so we may assume $p = 0$. For $q = 1$ this is
Theorem~\ref{several-complex-variables:theorem-dbar-poincare}. For general $q$ the same induction applies: if $\alpha$ involves only $d\bar z^1, \ldots, d\bar z^k$, write
$\alpha = d\bar z^k \wedge \alpha_1 + \alpha_2$ with $\alpha_1, \alpha_2$ not involving $d\bar z^k, \ldots, d\bar z^n$; the equation $\bar\partial\alpha = 0$ implies that the coefficients of $\alpha_1$ are holomorphic in
$z_{k+1}, \ldots, z_n$. Solving $\partial g/\partial\bar z_k = $ (coefficient of $\alpha_1$) with the Cauchy transform of Lemma~\ref{several-complex-variables:lemma-cauchy-transform}, after a cut-off in $z_k$
alone, gives a form $\gamma$ with $\alpha - \bar\partial\gamma$ involving only $d\bar z^1, \ldots, d\bar z^{k-1}$ on a slightly smaller polydisc; by induction on $k$ we conclude. The details are exactly as in the proof of
Theorem~\ref{several-complex-variables:theorem-dbar-poincare}; see \cite[Chapter 0, Section 2]{griffiths-harris}.
\end{proof}

\begin{theorem}[Dolbeault]\label{dolbeault:theorem-dolbeault}
The sequence of sheaves $0 \to \Omega^p_X(E) \to \mathcal{A}^{p,0}(E) \xrightarrow{\bar\partial_E} \mathcal{A}^{p,1}(E) \xrightarrow{\bar\partial_E} \cdots$ is a fine resolution of the sheaf
$\Omega^p_X(E) = \Omega^p_X \otimes E$ of holomorphic $E$-valued $p$-forms. Consequently
\[
H^{p,q}_{\bar\partial}(X, E) \cong H^q(X, \Omega^p_X \otimes E).
\]
\end{theorem}

\begin{proof}
Exactness at $\mathcal{A}^{p,0}$ says that the $E$-valued $(p,0)$-forms killed by $\bar\partial_E$ are the holomorphic ones (Proposition~\ref{complex-manifolds:proposition-dbar}), and exactness at
$\mathcal{A}^{p,q}$ for $q \geq 1$ holds on stalks by Lemma~\ref{dolbeault:lemma-dbar-poincare-pq}, applied in a holomorphic frame of $E$ on small polydiscs. The sheaves $\mathcal{A}^{p,q}(E)$ are modules over the
sheaf of smooth functions, hence fine and acyclic (Theorem~\ref{sheaf-cohomology:theorem-soft-acyclic}); apply Proposition~\ref{sheaf-cohomology:proposition-hypercohomology}(3).
\end{proof}

\begin{theorem}[Holomorphic Poincar\'e lemma]\label{dolbeault:theorem-holomorphic-de-rham}
The holomorphic de Rham complex $0 \to \CC_X \to \Omega^0_X \xrightarrow{\partial} \Omega^1_X \xrightarrow{\partial} \cdots \to \Omega^n_X \to 0$ is exact. Hence
$H^k(X; \CC) \cong \mathbb{H}^k(X, \Omega^\bullet_X)$, and there is a spectral sequence (the \emph{Fr\"olicher} or \emph{Hodge-to-de Rham} spectral sequence)
\[
E_1^{p,q} = H^q(X, \Omega^p_X) \Rightarrow H^{p+q}(X; \CC).
\]
\end{theorem}

\begin{proof}
On holomorphic forms $d = \partial$. Consider the double complex of sheaves $\mathcal{A}^{p,q}$ with differentials $\partial$ and $\bar\partial$, whose total complex is the complexified de Rham complex. By
Theorem~\ref{dolbeault:theorem-dolbeault} each column $(\mathcal{A}^{p,\bullet}, \bar\partial)$ is a resolution of $\Omega^p_X$; by Lemma~\ref{homological-algebra:lemma-acyclic-assembly}, applied to stalks, the
inclusion $\Omega^\bullet_X \to \mathrm{Tot}(\mathcal{A}^{\bullet,\bullet}) = \mathcal{A}^\bullet_\CC$ is a quasi-isomorphism. The de Rham complex is a resolution of $\CC_X$
(Theorem~\ref{de-rham:theorem-de-rham}), so $\Omega^\bullet_X$ has cohomology sheaf $\CC_X$ in degree $0$ and none elsewhere, which is the exactness. The isomorphism with $H^k(X; \CC)$ follows from
Proposition~\ref{sheaf-cohomology:proposition-hypercohomology}(1) and Theorem~\ref{de-rham:theorem-de-rham}, and the spectral sequence is the first one of
Proposition~\ref{sheaf-cohomology:proposition-hypercohomology}(2).
\end{proof}

\section{Harmonic theory for $\bar\partial$}\label{dolbeault:section-harmonic}

\begin{definition}\label{dolbeault:definition-dbar-laplacian}
Let $X$ be compact with a Hermitian metric and $E$ a Hermitian holomorphic bundle. With the $L^2$ inner products on $\mathcal{A}^{p,q}(X, E)$, let $\bar\partial_E^*$ be the formal adjoint of $\bar\partial_E$
(Definition~\ref{elliptic-operators:definition-formal-adjoint}) and $\Box_E = \bar\partial_E\bar\partial_E^* + \bar\partial_E^*\bar\partial_E$, the \emph{$\bar\partial$-Laplacian}. The space of
\emph{$\bar\partial$-harmonic forms} is $\mathcal{H}^{p,q}(X, E) = \ker\Box_E$.
\end{definition}

\begin{theorem}\label{dolbeault:theorem-dbar-hodge}
$\Box_E$ is elliptic, with principal symbol $\frac12|\xi|^2$. Consequently $\mathcal{H}^{p,q}(X, E)$ is finite-dimensional, there is an orthogonal decomposition
$\mathcal{A}^{p,q}(X, E) = \mathcal{H}^{p,q}(X, E) \oplus \bar\partial_E\mathcal{A}^{p,q-1}(X, E) \oplus \bar\partial_E^*\mathcal{A}^{p,q+1}(X, E)$, and every Dolbeault class has a unique harmonic
representative:
\[
H^q(X, \Omega^p_X \otimes E) \cong H^{p,q}_{\bar\partial}(X, E) \cong \mathcal{H}^{p,q}(X, E).
\]
In particular the Hodge numbers $h^{p,q}(X) = \dim H^q(X, \Omega^p_X)$ of a compact complex manifold are finite.
\end{theorem}

\begin{proof}
The symbol of $\bar\partial$ is $i\xi^{0,1} \wedge -$ and that of $\bar\partial^*$ is $-i\iota_{\xi^{0,1}}$, so the symbol of $\Box_E$ is $\xi^{0,1} \wedge \iota + \iota\,\xi^{0,1} \wedge = |\xi^{0,1}|^2 = \frac12|\xi|^2$, by the
same computation as in Exercise~\ref{elliptic-operators:exercise-symbol-laplacian}. The rest is the proof of Theorems~\ref{hodge-theorem:theorem-decomposition}
and~\ref{hodge-theorem:theorem-hodge} with $d$, $\delta$ replaced by $\bar\partial_E$, $\bar\partial_E^*$ (Remark~\ref{hodge-theorem:remark-elliptic-complexes}): the only inputs were ellipticity, self-adjointness and
$\bar\partial_E^2 = 0$. Combine with Theorem~\ref{dolbeault:theorem-dolbeault}.
\end{proof}

\begin{theorem}[Serre duality]\label{dolbeault:theorem-serre-duality}
Let $X$ be a compact complex manifold of dimension $n$ and $E$ a holomorphic vector bundle. The pairing
\[
H^q(X, E) \times H^{n-q}(X, K_X \otimes E^\vee) \to H^n(X, K_X) \xrightarrow{\ \int\ } \CC, \qquad ([\alpha], [\beta]) \mapsto \int_X\alpha \wedge \beta,
\]
computed with Dolbeault representatives of types $(0,q)$ and $(n, n-q)$ and the evaluation $E \otimes E^\vee \to \CC$, is perfect. In particular $h^{p,q} = h^{n-p,n-q}$ for the Hodge numbers.
\end{theorem}

\begin{proof}
The pairing is well defined by Stokes' theorem, since $\int_X\bar\partial(\alpha \wedge \gamma) = \int_Xd(\alpha \wedge \gamma) = 0$ for forms of total type $(n, n-1)$. Choose Hermitian metrics and define the conjugate-linear
operator $\bar*_E \colon \mathcal{A}^{0,q}(E) \to \mathcal{A}^{n,n-q}(E^\vee)$ by $\alpha \wedge \bar*_E\beta = \langle\alpha, \beta\rangle\,\mathrm{vol}$, using the metric to identify $\overline{E}$ with $E^\vee$. As
for the real Hodge star, $\bar\partial^* = -\bar*\bar\partial\bar*$, so $\bar*_E$ maps $\Box_E$-harmonic forms bijectively onto $\Box_{K \otimes E^\vee}$-harmonic forms. For a nonzero harmonic $\alpha$,
$\int_X\alpha \wedge \bar*_E\alpha = \|\alpha\|^2 > 0$, so the pairing is nondegenerate on harmonic representatives, which by Theorem~\ref{dolbeault:theorem-dbar-hodge} represent all classes. For the Hodge
numbers take $E = \Omega^p_X$, so $K_X \otimes E^\vee \cong \Omega^{n-p}_X$ by the perfect pairing $\Omega^p \otimes \Omega^{n-p} \to K_X$.
\end{proof}

\section{Examples and the Fr\"olicher inequality}\label{dolbeault:section-examples}

\begin{proposition}\label{dolbeault:proposition-frolicher-inequality}
For a compact complex manifold, $b_k(X) \leq \sum_{p+q=k}h^{p,q}(X)$, and $\chi(X) = \sum_{p,q}(-1)^{p+q}h^{p,q}(X)$. Equality holds for all $k$ if and only if the Fr\"olicher spectral sequence degenerates at
$E_1$.
\end{proposition}

\begin{proof}
In a spectral sequence of finite-dimensional vector spaces, $\dim E_{r+1}^{p,q} \leq \dim E_r^{p,q}$ and $\sum_{p+q=k}\dim E_\infty^{p,q} = b_k$ (Theorem~\ref{spectral-sequences:theorem-filtered-complex}), so
$b_k \leq \sum_{p+q=k}h^{p,q}$, with equality for all $k$ if and only if all differentials $d_r$, $r \geq 1$, vanish. The Euler characteristic is unchanged when passing to cohomology of a differential of
total degree one, so it is the same on every page.
\end{proof}

\begin{example}\label{dolbeault:example-hodge-numbers}
\begin{enumerate}
\item For a complex torus $X = \CC^n/\Lambda$, the forms $dz^I \wedge d\bar z^J$ with constant coefficients are harmonic for the flat metric and span each $\mathcal{H}^{p,q}$, so
$h^{p,q} = \binom{n}{p}\binom{n}{q}$, and $\sum_{p+q=k}h^{p,q} = \binom{2n}{k} = b_k$: the Fr\"olicher spectral sequence degenerates.
\item For $\CC\PP^n$, $h^{p,q} = 1$ if $p = q \leq n$ and $0$ otherwise. Since $b_{2k} = 1$ and $b_{2k+1} = 0$, this follows from the Hodge decomposition (Chapter~\ref{hodge-decomposition}) once one knows
$h^{k,k} \geq 1$, which is witnessed by the class of $\omega_{FS}^k$; see \cite[Chapter 0, Section 7]{griffiths-harris}.
\item For the Hopf surface of Exercise~\ref{complex-manifolds:exercise-hopf-surface}, $b_1 = 1$, $h^{0,1} = 1$ and $h^{1,0} = 0$. So $b_1 = h^{1,0} + h^{0,1}$ (the Fr\"olicher spectral sequence of a compact
complex surface always degenerates), but $h^{1,0} \neq h^{0,1}$: the Hodge symmetry $h^{p,q} = h^{q,p}$ of compact K\"ahler manifolds fails \cite[Chapter IV]{bhpv}.
\end{enumerate}
\end{example}

\begin{remark}\label{dolbeault:remark-vanishing}
Cohomology of line bundles on projective space is computed by Bott's formula: $H^q(\CC\PP^n, \mathcal{O}(k))$ vanishes for $0 < q < n$ and all $k$, $H^0 = $ homogeneous polynomials of degree $k$, and
$H^n(\CC\PP^n, \mathcal{O}(k)) \cong H^0(\CC\PP^n, \mathcal{O}(-k-n-1))^\vee$ by Serre duality, since $K_{\CC\PP^n} = \mathcal{O}(-n-1)$ (Proposition~\ref{holomorphic-bundles:proposition-adjunction}).
Vanishing theorems for positive bundles are the subject of Chapter~\ref{kodaira}; the algebraic computation is in Chapter~\ref{coherent-cohomology}.
\end{remark}

\section{Exercises}\label{dolbeault:section-exercises}

\begin{exercise}\label{dolbeault:exercise-riemann-surface}
Let $X$ be a compact Riemann surface of genus $g$. Using Serre duality and the Fr\"olicher inequality, show that $h^{0,0} = h^{1,1} = 1$ and $h^{1,0} = h^{0,1}$, and that $h^{1,0} + h^{0,1} \geq 2g$ with equality
if and only if the Fr\"olicher spectral sequence degenerates (it always does in dimension one, by Chapter~\ref{hodge-decomposition}).
\end{exercise}

\begin{solution}
$h^{0,0} = 1$ since holomorphic functions are constant (Proposition~\ref{complex-manifolds:proposition-no-functions}); Serre duality gives $h^{1,1} = h^{0,0}$ and $h^{1,0} = h^{0,1}$. The inequality is
Proposition~\ref{dolbeault:proposition-frolicher-inequality} for $k = 1$, with $b_1 = 2g$.
\end{solution}

\begin{exercise}\label{dolbeault:exercise-dolbeault-plane}
Show that $H^{0,1}_{\bar\partial}(\CC) = 0$, but that $H^{0,1}_{\bar\partial}(\CC^2 \setminus \{0\})$ is infinite-dimensional.
\end{exercise}

\begin{exercise}\label{dolbeault:exercise-serre-duality-curve}
For a compact Riemann surface $X$ and a line bundle $L$, show that $H^1(X, L) \cong H^0(X, K_X \otimes L^\vee)^\vee$, and deduce that $H^1(X, L) = 0$ if $\deg L > 2g - 2$.
\end{exercise}

\begin{exercise}\label{dolbeault:exercise-euler-characteristic-bundle}
Show that the holomorphic Euler characteristic $\chi(X, E) = \sum_q(-1)^q\dim H^q(X, E)$ of a compact complex manifold is the index of the elliptic operator
$\bar\partial_E + \bar\partial_E^* \colon \mathcal{A}^{0,\mathrm{even}}(E) \to \mathcal{A}^{0,\mathrm{odd}}(E)$.
\end{exercise}
